% ============================================================
% Round 06: DDPO / Diffusion RL
% ============================================================

\subsection{Round 06：DDPO 与 Diffusion RL}

\subsubsection{Highlight}

\begin{conceptbox}
\textbf{本章相关脉络：}
\begin{itemize}[leftmargin=1.5em]
  \item \textbf{DDPO}：arXiv 2023-05-22，Training Diffusion Models with Reinforcement Learning，把 denoising process 写成多步 MDP，并用 policy gradient / PPO 类方法优化最终 reward。论文：\url{https://arxiv.org/abs/2305.13301}，代码：\url{https://github.com/kvablack/ddpo-pytorch}
  \item \textbf{DPOK}：arXiv 2023-05-25，NeurIPS 2023，Reinforcement Learning for Fine-tuning Text-to-Image Diffusion Models，用 online RL + KL regularization 微调 text-to-image diffusion model。论文：\url{https://arxiv.org/abs/2305.16381}，代码：\url{https://github.com/google-research/google-research/tree/master/dpok}
  \item \textbf{TRL / Diffusers DDPOTrainer}：Hugging Face 后续把 DDPO 训练接入 TRL + Diffusers，用于基于 reward function 微调 Stable Diffusion。文档：\url{https://huggingface.co/docs/diffusers/en/training/ddpo}
\end{itemize}
\end{conceptbox}

\subsubsection{本章要解决的问题}

上一章 Diffusion-DPO 用的是离线偏好对：

\[
(I_\text{src},c,I_w,I_l)
\]

它的目标是：

\[
\text{相对 rejected，更偏向 chosen}
\]

这一章进入 online RL。我们不只拿已有的 winner / loser，而是让当前模型自己采样编辑结果：

\[
I_\text{edit}
\sim
\pi_\theta(\cdot\mid I_\text{src},c)
\]

然后用 reward model / scorer 打分：

\[
R
=
R(I_\text{src},c,I_\text{edit})
\]

再用这个 reward 更新模型。

\begin{intubox}
\textbf{直觉：} DPO 像是在做离线选择题：“这两个结果哪个更好”；DDPO / Diffusion RL 像是在做在线练习：“你自己生成一批图，得到分数，然后把高分轨迹变得更可能”。
\end{intubox}

\subsubsection{Diffusion 采样为什么能看成 MDP}

Diffusion 采样是一条去噪轨迹：

\[
z_T
\to
z_{T-1}
\to
\cdots
\to
z_0
\]

其中 $z_T$ 是噪声，$z_0$ 解码后得到最终图像：

\[
I_\text{edit}
=
D_\text{VAE}(z_0)
\]

在图像编辑里，条件是：

\[
x=(I_\text{src},c)
\]

所以每一步反向转移是：

\[
p_\theta(z_{t-1}\mid z_t,x,t)
\]

把它翻译成 RL 的 MDP 语言：

\begin{center}
{\small
\begin{tabular}{p{3cm}p{4.6cm}p{5cm}}
\hline
\textbf{RL 概念} & \textbf{Diffusion 对应物} & \textbf{说明} \\
\hline
State
&
$s_t=(z_t,t,x)$
&
当前 noisy latent、时间步、编辑条件
\\
Action
&
$a_t=z_{t-1}$
&
模型选择下一步 denoised latent
\\
Policy
&
$\pi_\theta(a_t\mid s_t)=p_\theta(z_{t-1}\mid z_t,x,t)$
&
扩散模型的一步反向转移
\\
Trajectory
&
$\tau=(z_T,z_{T-1},\ldots,z_0)$
&
完整去噪路径
\\
Reward
&
$R(I_\text{src},c,D_\text{VAE}(z_0))$
&
最终图像得分，通常是 terminal reward
\\
\hline
\end{tabular}
}
\end{center}

\begin{summarybox}
\textbf{最重要的对应关系：}\\
文本 RL 里 action 是 token；diffusion RL 里 action 是每一步 denoising transition，也就是从 $z_t$ 走到 $z_{t-1}$。
\end{summarybox}

\subsubsection{每一步 Policy 的 Logprob}

为了做 PPO / DDPO，我们需要每一步 action 的 logprob：

\[
\log
p_\theta(z_{t-1}\mid z_t,x,t)
\]

在很多 diffusion scheduler 里，一步反向转移可以写成高斯：

\[
p_\theta(z_{t-1}\mid z_t,x,t)
=
\mathcal{N}
\left(
z_{t-1};
\mu_\theta(z_t,t,x),
\sigma_t^2 I
\right)
\]

其中：

\begin{itemize}[leftmargin=1.5em]
  \item $\mu_\theta(z_t,t,x)$ 是模型根据当前 noisy latent 预测出来的下一步均值；
  \item $\sigma_t$ 是 scheduler 给定的随机噪声尺度；
  \item $z_{t-1}$ 是这次真实采样出来的下一步 latent。
\end{itemize}

因此 logprob 是：

\[
\log
p_\theta(z_{t-1}\mid z_t,x,t)
=
-
\frac{
\left\|
z_{t-1}
-
\mu_\theta(z_t,t,x)
\right\|^2
}{
2\sigma_t^2
}
-
\frac{d}{2}\log(2\pi\sigma_t^2)
\]

其中 $d$ 是 latent 维度。

\begin{intubox}
\textbf{直觉：} 如果这次采样出的 $z_{t-1}$ 很靠近模型均值 $\mu_\theta$，那这一步 logprob 高；如果偏得很远，logprob 低。
\end{intubox}

\subsubsection{整条 Trajectory 的 Logprob}

一条完整去噪轨迹的概率是每一步概率相乘：

\[
p_\theta(\tau\mid x)
=
p(z_T)
\prod_{t=1}^{T}
p_\theta(z_{t-1}\mid z_t,x,t)
\]

取 log 就变成求和：

\[
\log p_\theta(\tau\mid x)
=
\log p(z_T)
+
\sum_{t=1}^{T}
\log
p_\theta(z_{t-1}\mid z_t,x,t)
\]

因为初始噪声 $p(z_T)$ 通常和 $\theta$ 无关，所以更新模型时主要看：

\[
\sum_{t=1}^{T}
\log
p_\theta(z_{t-1}\mid z_t,x,t)
\]

这和文本侧非常像：

\[
\log \pi_\theta(y\mid x)
=
\sum_i
\log \pi_\theta(y_i\mid x,y_{<i})
\]

只是文本求和的是 token logprob，diffusion 求和的是 denoising transition logprob。

\begin{summarybox}
\textbf{类比：}\\
文本序列：
\[
y_1,y_2,\ldots,y_n
\]
的 logprob 是 token logprob 之和。\\
Diffusion trajectory：
\[
z_T,z_{T-1},\ldots,z_0
\]
的 logprob 是 denoising step logprob 之和。
\end{summarybox}

\subsubsection{最终 Reward 如何更新中间步骤}

现在模型采样出一条轨迹：

\[
\tau=(z_T,z_{T-1},\ldots,z_0)
\]

解码最终图像：

\[
I_\text{edit}=D_\text{VAE}(z_0)
\]

然后 reward model 给一个分数：

\[
R(\tau)
=
R(I_\text{src},c,I_\text{edit})
\]

目标是最大化期望 reward：

\[
J(\theta)
=
\mathbb{E}_{\tau\sim p_\theta(\tau\mid x)}
\left[
R(\tau)
\right]
\]

用 REINFORCE / policy gradient：

\[
\nabla_\theta J(\theta)
=
\mathbb{E}
\left[
R(\tau)
\nabla_\theta
\log p_\theta(\tau\mid x)
\right]
\]

把 trajectory logprob 展开：

\[
\nabla_\theta J(\theta)
=
\mathbb{E}
\left[
R(\tau)
\sum_{t=1}^{T}
\nabla_\theta
\log p_\theta(z_{t-1}\mid z_t,x,t)
\right]
\]

通常会减去 baseline，得到 advantage：

\[
A(\tau)
=
R(\tau)-b(x)
\]

于是：

\[
\nabla_\theta J(\theta)
=
\mathbb{E}
\left[
A(\tau)
\sum_{t=1}^{T}
\nabla_\theta
\log p_\theta(z_{t-1}\mid z_t,x,t)
\right]
\]

\begin{intubox}
\textbf{直觉：} 如果最终图像得分高，就提高这条去噪轨迹中每一步 transition 的概率；如果最终图像得分低，就降低这条轨迹中每一步 transition 的概率。
\end{intubox}

这就是“最终 reward 如何回传到中间步骤”的最小答案：

\[
\text{不是直接知道哪一步错了}
\]

\[
\text{而是用同一个最终 advantage 给整条轨迹的每一步 logprob 加权}
\]

\begin{warnbox}
\textbf{信用分配问题：} 这种 terminal reward 很粗糙。图像坏了可能只因为某几个步骤，但 policy gradient 会把同一个 reward 信号分给很多 denoising steps，所以方差高、训练贵，也容易不稳定。
\end{warnbox}

\subsubsection{DDPO：把 PPO Ratio 用到 Denoising Step}

实际训练通常不是每采样一条轨迹就立刻更新一次，而是先用旧模型采样：

\[
\tau\sim p_{\theta_\text{old}}(\tau\mid x)
\]

记录每一步旧 logprob：

\[
\log p_{\theta_\text{old}}(z_{t-1}\mid z_t,x,t)
\]

然后更新新模型时，重新计算：

\[
\log p_\theta(z_{t-1}\mid z_t,x,t)
\]

每一步 importance ratio 是：

\[
r_t(\theta)
=
\exp
\left(
\log p_\theta(z_{t-1}\mid z_t,x,t)
-
\log p_{\theta_\text{old}}(z_{t-1}\mid z_t,x,t)
\right)
\]

这和文本 PPO 里的 token ratio 完全同构：

\[
r_i(\theta)
=
\frac{
\pi_\theta(y_i\mid x,y_{<i})
}{
\pi_{\theta_\text{old}}(y_i\mid x,y_{<i})
}
\]

只不过 diffusion 里比较的是：

\[
\text{同一个 denoising transition 在新旧模型下的概率}
\]

PPO-style clipped objective 可以写成：

\[
\mathcal{L}_\text{DDPO}
=
-
\mathbb{E}
\left[
\sum_{t=1}^{T}
\min
\left(
r_t(\theta)A_t,
\text{clip}(r_t(\theta),1-\epsilon,1+\epsilon)A_t
\right)
\right]
\]

如果只有最终 reward，最简单可以让每一步共享同一个 advantage：

\[
A_t=A(\tau)
\]

也可以按时间步做更细的 advantage / discount / normalization。

\begin{summarybox}
\textbf{DDPO 一句话：}\\
DDPO 把 diffusion 每一步反向转移当成 policy action，记录旧模型采样这一步的 logprob，再用 PPO ratio 更新新模型，让高 reward 的 denoising trajectory 更可能出现。
\end{summarybox}

\subsubsection{图像编辑版 DDPO}

Text-to-image DDPO 的条件是文本 prompt：

\[
x=c_\text{text}
\]

图像编辑 DDPO 的条件是：

\[
x=(I_\text{src},c)
\]

每一步转移写成：

\[
p_\theta(z_{t-1}\mid z_t,I_\text{src},c,t)
\]

最终 reward 写成：

\[
R
=
R(I_\text{src},c,D_\text{VAE}(z_0))
\]

所以图像编辑版 DDPO 的完整闭环是：

\[
(I_\text{src},c)
\Rightarrow
z_T
\Rightarrow
z_{T-1}
\Rightarrow
\cdots
\Rightarrow
z_0
\Rightarrow
I_\text{edit}
\Rightarrow
R
\Rightarrow
\text{update denoising logprobs}
\]

这里的 reward 通常需要综合：

\[
R
=
w_\text{inst}R_\text{inst}
+
w_\text{pres}R_\text{pres}
+
w_\text{qual}R_\text{qual}
+
w_\text{pref}R_\text{pref}
\]

\begin{warnbox}
\textbf{图像编辑更难：} Text-to-image 主要看 prompt-image alignment 和质量；image editing 还必须同时看是否保留原图。不加 preservation reward，RL 很容易把编辑任务优化成“重新生成一张高分图”。
\end{warnbox}

\subsubsection{DPOK：为什么还要 KL Regularization}

RL 直接最大化 reward 很容易把模型推离原始分布。为了稳定，DPOK 这类方法会显式加入 KL regularization：

\[
\max_\theta
\mathbb{E}_{\tau\sim p_\theta}
\left[
R(\tau)
\right]
-
\lambda
\mathrm{KL}
\left(
p_\theta(\tau\mid x)
\;\|\;
p_\text{ref}(\tau\mid x)
\right)
\]

其中 $p_\text{ref}$ 是冻结的 reference diffusion model。

轨迹 KL 可以拆到每一步：

\[
\mathrm{KL}
\left(
p_\theta(\tau\mid x)
\;\|\;
p_\text{ref}(\tau\mid x)
\right)
\approx
\sum_{t=1}^{T}
\mathrm{KL}
\left(
p_\theta(z_{t-1}\mid z_t,x,t)
\;\|\;
p_\text{ref}(z_{t-1}\mid z_t,x,t)
\right)
\]

也可以在采样到的 transition 上用 logprob 差近似 KL penalty：

\[
r_\text{KL,t}
=
-
\lambda
\left[
\log p_\theta(z_{t-1}\mid z_t,x,t)
-
\log p_\text{ref}(z_{t-1}\mid z_t,x,t)
\right]
\]

这样总 reward 可以写成：

\[
R_\text{total}
=
R_\text{external}
+
\sum_{t=1}^{T}
r_\text{KL,t}
\]

\begin{intubox}
\textbf{直觉：} 外部 reward 负责“往高分方向走”；KL 负责“不要离原模型太远”。这和文本 RLHF 里 PPO + reference KL 的作用一致。
\end{intubox}

\subsubsection{DDPO、DPOK、Diffusion-DPO 的区别}

\begin{center}
{\small
\begin{tabular}{p{2.7cm}p{3.8cm}p{4cm}p{3.2cm}}
\hline
\textbf{方法} & \textbf{数据来源} & \textbf{优化信号} & \textbf{特点} \\
\hline
Diffusion-DPO
&
离线偏好对 $(I_w,I_l)$
&
chosen / rejected 的相对 denoising error
&
稳定，依赖偏好数据覆盖
\\
DDPO
&
模型在线采样轨迹
&
最终 reward 加权每步 logprob
&
能探索，但方差高、采样贵
\\
DPOK
&
模型在线采样轨迹
&
reward + KL regularization
&
更强调不偏离 reference
\\
\hline
\end{tabular}
}
\end{center}

\subsubsection{为什么不直接对 Reward 反向传播}

有些 reward 是可微的，例如某些 CLIP similarity：

\[
R_\text{CLIP}(I,c)
\]

理论上可以从 reward 对图像求梯度，再穿过 VAE 和 denoising process 反传。

但很多实际 reward 不好直接反传：

\begin{itemize}[leftmargin=1.5em]
  \item VLM judge 输出离散评分；
  \item 人类偏好 reward model 可能是黑盒 API；
  \item JPEG compressibility、OCR、object counting 等指标可能不可微或不稳定；
  \item 完整反传穿过几十步 denoising 很吃显存；
  \item 多 reward 组合时，有些项可微，有些项不可微。
\end{itemize}

Policy gradient 的优点是：

\[
\text{只要能给最终样本打分，就可以优化}
\]

它不要求 reward 对图像可微。

\begin{summarybox}
\textbf{DDPO 的工程价值：}\\
DDPO 把 diffusion model 当成随机策略。只要能采样、能计算每步 logprob、能给最终图像打分，就可以用 policy gradient 优化，即使 reward 本身不可微。
\end{summarybox}

\subsubsection{一个最小 DDPO 训练流程}

\begin{enumerate}[leftmargin=1.5em]
  \item 取一批编辑任务 $(I_\text{src},c)$；
  \item 用当前 old policy 采样多条 denoising trajectories；
  \item 解码每条轨迹的最终图像 $I_\text{edit}$；
  \item 用 reward model / scorer 计算 $R(I_\text{src},c,I_\text{edit})$；
  \item 对同一批样本做 reward normalization，得到 advantage $A$；
  \item 记录旧 logprob $\log p_{\theta_\text{old}}(z_{t-1}\mid z_t,x,t)$；
  \item 用当前模型重新计算新 logprob；
  \item 构造每步 ratio $r_t(\theta)$；
  \item 用 PPO clipped objective 更新模型；
  \item 周期性刷新 old policy，继续采样。
\end{enumerate}

压缩成一行：

\[
\text{sample trajectory}
\Rightarrow
\text{score final image}
\Rightarrow
\text{compute advantage}
\Rightarrow
\text{update step logprobs}
\]

\subsubsection{和文本 PPO / GRPO 的对应关系}

\begin{center}
{\small
\begin{tabular}{p{3cm}p{4.8cm}p{4.8cm}}
\hline
\textbf{概念} & \textbf{文本 PPO / GRPO} & \textbf{DDPO / Diffusion RL} \\
\hline
Policy
&
$\pi_\theta(y_i\mid x,y_{<i})$
&
$p_\theta(z_{t-1}\mid z_t,x,t)$
\\
Action
&
下一个 token
&
下一步 latent $z_{t-1}$
\\
Trajectory
&
完整回答 $y$
&
完整去噪路径 $\tau$
\\
Logprob
&
token logprob 求和
&
transition logprob 求和
\\
Reward
&
回答级 reward
&
最终图像 reward
\\
Ratio
&
新旧 token probability ratio
&
新旧 denoising transition probability ratio
\\
\hline
\end{tabular}
}
\end{center}

\begin{intubox}
\textbf{记忆法：} 文本 PPO 是“奖励一段回答里的 token 轨迹”；DDPO 是“奖励一张图背后的 denoising 轨迹”。
\end{intubox}

\subsubsection{DDPO 的主要问题}

DDPO 很直接，但问题也明显：

\begin{itemize}[leftmargin=1.5em]
  \item \textbf{采样贵：} 每次更新前都要跑完整 diffusion sampling；
  \item \textbf{方差高：} 最终 reward 要分给很多 denoising steps；
  \item \textbf{reward hacking：} 模型可能钻 reward model 的空子；
  \item \textbf{scheduler 依赖：} logprob 的定义和采样 scheduler 强相关；
  \item \textbf{KL 控制困难：} KL 太弱会漂，太强又学不动；
  \item \textbf{编辑保持脆弱：} 图像编辑里如果 reward 没有约束原图保持，模型会过度编辑。
\end{itemize}

\begin{warnbox}
\textbf{为什么后面还要学 Flow-GRPO / HP-Edit：}\\
DDPO 给出了 diffusion RL 的基础框架，但现代图像编辑模型可能是 flow-based，也可能需要更细的 task-aware reward、group-relative advantage、dense evaluator 和偏好数据闭环。HP-Edit 这类工作就是在这些问题上继续往前走。
\end{warnbox}

\subsubsection{本章最小结论}

\begin{summarybox}
\textbf{DDPO 的核心：}\\
把 diffusion denoising process 看成多步 MDP：$z_t$ 是 state，$z_{t-1}$ 是 action，每一步反向转移有 logprob，整条轨迹的 logprob 是所有 step logprob 之和。最终图像 reward 通过 policy gradient 加权这些 step logprob，从而让高 reward 的去噪轨迹更可能出现。图像编辑版 DDPO 只是把条件从 text prompt 换成 $(I_\text{src},c)$，同时 reward 必须额外约束编辑正确性和原图保持。
\end{summarybox}

\subsubsection{本轮检查题}

\begin{enumerate}[leftmargin=1.5em]
  \item 为什么 diffusion denoising process 可以看成一个 MDP？
  \item 在 DDPO 里，state、action、policy、trajectory、reward 分别是什么？
  \item 整条 denoising trajectory 的 logprob 为什么是每一步 logprob 的和？
  \item 最终图像 reward 是如何作用到中间 denoising steps 的？
  \item DDPO 的 importance ratio 和文本 PPO 的 ratio 有什么对应关系？
  \item 图像编辑版 DDPO 为什么必须加入 preservation reward？
  \item DPO、DDPO、DPOK 三者最主要的区别是什么？
\end{enumerate}
